%% ====================================================================
%%  APPENDIX A: FULL PROOFS
%% ====================================================================
\appendix
\section{Full Proofs}
\label{app:proofs}

\subsection{Graph-Depth Model (Theorem~\ref{thm:optimal-depth})}
\label{app:graph-depth}

\begin{proof}
\noindent\textbf{Setup.}
The builder's sub-problem for model~$M$ is
$\min_{D \ge 1} L(M,D) + \mu \cdot D \cdot \kappa(M)$,
where $L(M,D) = \sum_s \pi_s (-\log(1 - r_s^D))$ and $\mu > 0$ is the shadow price of the cost budget.

\noindent\textbf{Part (a): First-order condition.}
Differentiating $L$ with respect to~$D$:
\[
  \frac{\partial L}{\partial D}
  = \sum_s \pi_s \cdot \frac{-r_s^D \log r_s}{1 - r_s^D}
  = -F(D; M).
\]
The FOC $\partial L/\partial D + \mu \kappa = 0$ gives $F(\Dstar; M) = \mu \cdot \kappa(M)$.

\noindent\textbf{Part (b): Uniqueness.}
Define $f_s(D) = r_s^D (-\log r_s) / (1 - r_s^D)$.
Using the quotient rule with $u = r_s^D$, $v = -\log r_s$ (constant in~$D$), $w = 1 - r_s^D$:
\[
  \frac{df_s}{dD} = \frac{(-\log r_s) \cdot r_s^D \log r_s}{(1-r_s^D)^2}
  = \frac{-r_s^D (\log r_s)^2}{(1-r_s^D)^2} < 0.
\]
Since each $f_s$ is strictly decreasing, so is $F = \sum_s \pi_s f_s$.
As $D \to \infty$, $F \to 0$; as $D \to 0^+$, $F \to \infty$.
For any $\mu \kappa > 0$, the equation $F(\Dstar) = \mu \kappa$ has a unique solution.

\noindent\textbf{Part (c): Monotonicity in competence.}
Apply the implicit function theorem to $\Phi(\Dstar, t_0) = F(\Dstar; M) - \mu\kappa = 0$.
We have $\partial \Phi / \partial D = dF/dD < 0$ (from part~b).

We need $\partial F / \partial t_0$.
Under monotone competence, $p_s$ increases in~$t_0$ and $r_s = 1-p_s$ decreases.
Define $g(r) = r^D(-\log r)/(1-r^D)$.
Then:
\[
  \frac{dg}{dr} = \frac{r^{D-1}[D(-\log r) - 1 + r^D]}{(1-r^D)^2}.
\]
Define $h(r) = -D\log r - 1 + r^D$.
Then $h(1) = 0$ and $h'(r) = -D/r + Dr^{D-1} = D(r^D - 1)/r < 0$ for $r \in (0,1)$.
Since $h$ is decreasing with $h(1) = 0$, we have $h(r) > 0$ for all $r \in (0,1)$, giving $dg/dr > 0$.

By the IFT:
$d\Dstar/dp_s = -(\partial F/\partial p_s) / (\partial F/\partial D) < 0$.
Under monotone competence ($dp_s/dt_0 > 0$):
$d\Dstar/dt_0 = \sum_s (d\Dstar/dp_s)(dp_s/dt_0) < 0$.

\noindent\textbf{Part (d): Monotonicity in cost ratio.}
From $F(\Dstar) = \mu\kappa$: differentiating gives $(dF/dD)(d\Dstar/d\kappa) = \mu$.
Since $dF/dD < 0$: $d\Dstar/d\kappa < 0$.
Thus $\Dstar$ is decreasing in~$\kappa(M)$, hence increasing in $\kappabig/\kappa(M)$.
\end{proof}

\subsection{Crossover Theorem (Theorem~\ref{thm:crossover})}
\label{app:crossover}

\begin{proof}
\noindent\textbf{Part (a): Existence.}
$L(M_s, D)$ is continuous and strictly decreasing in~$D$, with $L(M_s, D) \to 0$ as $D \to \infty$.
Since $L(M_l, 1) > 0$, by the intermediate value theorem there exists $\Dcross$ with $L(M_s, \Dcross) = L(M_l, 1)$.
Uniqueness follows from strict monotonicity.

\noindent\textbf{Part (b): Single mode.}
With $\Mn = 1$, $\pi_1 = 1$:
$-\log(1 - r_s^D) = -\log p_l$ gives $r_s^D = r_l$, so $\Dcross = \log r_l / \log r_s$.

\noindent\textbf{Part (c): Multi-mode.}
The LHS of $\sum_s \pi_s(-\log(1-r_s(M_s)^D)) = \sum_s \pi_s(-\log p(s,M_l))$ is strictly decreasing in~$D$, approaching $0$ as $D \to \infty$.
At $D = 1$: the LHS term for mode~$s$ is $-\log(1 - r_s^1) = -\log(1-(1-p_s)) = -\log p_s > -\log p_l$ (the RHS term), since $p_s < p_l$.
By the IVT, a unique solution exists.

\noindent\textbf{Part (d): Hardest-mode approximation.}
For hard tasks, $-\log(1-r_s^D) \approx r_s^D/(1-r_s^D)$, dominated by the mode~$s^*$ with largest~$r_s$.
The sum reduces approximately to the single-mode equation for $s^*$, giving $\Dcross \approx \max_s \log r_l(s)/\log r_s(s)$.
The approximation is an overestimate with mean ratio $1.8\times$ in numerical experiments.

\noindent\textbf{Part (e): Cost-adjusted dominance.}
For the small model to dominate with $\Dcross$ attempts vs.\ the large model with 1~attempt at cost~$\kappa_l$, we need $\Dcross \cdot \kappa_s \le \kappa_l$, i.e., $\Dcross \le R$.
\end{proof}

\subsection{Sample Complexity (Proposition~\ref{prop:sample-complexity})}
\label{app:sample-complexity}

\begin{proof}
\noindent\textbf{Step 1: Per-mode concentration.}
By Hoeffding's inequality, $\Prb(|\hat{p}_s - p_s| > \epsilon) \le 2\exp(-2n_s\epsilon^2)$.
Setting the RHS to $\dconf/\Mn$ and applying a union bound: with probability $\ge 1-\dconf$, simultaneously $|\hat{p}_s - p_s| \le \epsilon$ for all~$s$.

\noindent\textbf{Step 2: Lipschitz propagation.}
The competence term involves $\log p_s$ with derivative $1/p_s$, so its Lipschitz constant is $1/\pmin$.
For the information-gain term $G = \sum_s \pi_s \log(p_s / \bar{p})$ where $\bar{p} = \sum_s \pi_s p_s$: differentiating with respect to $p_s$ gives $\partial G / \partial p_s = \pi_s(1/p_s - 1/\bar{p})$.
Since $p_s \ge \pmin$ and $\bar{p} \le 1$: $|\partial G / \partial p_s| \le \pi_s / \pmin$.
Summing over modes: $\|\nabla_p G\|_1 = \sum_s |\partial G / \partial p_s| \le 1/\pmin$.
The routing mismatch $\mismatch = \sum_s \pi_s \log(q_s^*/q_s)$ has an analogous bound via the chain rule through $q_s^* = \pi_s p_s / \bar{p}$.
To achieve $\dacc$-accuracy in each term: set $\epsilon/\pmin = \dacc$, giving $\epsilon = \dacc\pmin$.

\noindent\textbf{Step 3: Combine.}
$n_0 = \log(2\Mn/\dconf) / (2\dacc^2\pmin^2)$.
Total: $N = \Mn \cdot n_0 = O(\Mn\log(\Mn/\dconf)/(\dacc^2\pmin^2))$.
\end{proof}

\subsection{Bernstein Corollary (Corollary~\ref{cor:bernstein})}
\label{app:bernstein}

\begin{proof}
Replace Hoeffding with Bernstein's inequality for Bernoulli random variables:
$\Prb(|\hat{p} - p| \ge \epsilon) \le 2\exp(-n\epsilon^2 / (2\sigma^2 + \frac{2}{3}\epsilon))$
where $\sigma^2 = p(1-p) \le p$ for $p \le 1/2$.

Setting $\epsilon = \dacc\pmin$ (from the Lipschitz step, unchanged):
\begin{align}
  n_0^{\mathrm{Bern}} &\ge \frac{(2\pmin + \frac{2}{3}\dacc\pmin)\log(2/\delta_s)}{(\dacc\pmin)^2}
  = \frac{(2 + \frac{2}{3}\dacc)\log(2/\delta_s)}{\dacc^2\pmin}.
\end{align}
For $\dacc \le 1$: the constant is at most $8/3 < 3$.
With $\delta_s = \dconf/\Mn$:
$n_0^{\mathrm{Bern}} = 3\log(2\Mn/\dconf)/(\dacc^2\pmin)$.

The improvement factor over Hoeffding is $n_0^{\mathrm{Hoeff}}/n_0^{\mathrm{Bern}} = 1/(6\pmin)$.
Bernstein dominates when $\pmin < 1/6 \approx 0.17$.
\end{proof}

\subsection{Lagrangian Non-Connection (Proposition~\ref{prop:lagrangian})}
\label{app:lagrangian}

\begin{proof}
\noindent\textbf{Step 1: KKT conditions for the $2 \times 2$ case.}
Fix model~$M$.
The Lagrangian for the routing sub-problem is:
$\mathcal{L}(q,\lambda,\nu) = \sum_s \pi_s(-\log q_s) + \lambda(\kappa(M)-B) + \nu(\sum_s q_s - 1)$.
KKT: $-\pi_s/q_s + \nu = 0$, giving $q_s^* = \pi_s/\nu$.
The simplex constraint gives $\nu = 1$, so $q_s^* = \pi_s$---the posterior-matched allocation.
The optimal objective is $\Hent(\pi)$, independent of~$M$.

\noindent\textbf{Step 2: The shadow price.}
There is one constraint (cost~$\le B$), hence one multiplier~$\lambda^*$.
If $\kappa(M) < B$: $\lambda^* = 0$.
If $\kappa(M) = B$: $\lambda^* \ge 0$, a single scalar.

\noindent\textbf{Step 3: Comparison.}
The four-term identity $\Delta = \log(\kappa_0/\kappa) + \varphi + G - \mismatch$ holds at every $(M,q)$---a pointwise algebraic identity.
The KKT conditions hold only at $(M^*,q^*)$---first-order optimality conditions.
For the four terms to be shadow prices, the problem would need four constraints whose sensitivities equal these quantities.
Standard operational constraints (cost, latency, memory) are aggregates that do not encode mode-level information structure.
Relaxing the cost budget changes the optimal model, affecting $\kappa$, $\varphi$, and $G$ simultaneously---not one term in isolation.

\noindent\textbf{Structural obstruction.}
Shadow prices measure $-\partial V^*/\partial b_k$ (sensitivity of the optimum to constraint relaxation).
The four terms are $\Delta_j(M^*,q^*)$ (components of the objective at the optimum).
These coincide only when relaxing constraint~$k$ shifts the optimum in a direction that changes exactly one term while leaving the others fixed---a generically false condition for operational constraints.
\end{proof}

\subsection{Proxy Validity (Theorem~\ref{thm:proxy-validity})}
\label{app:proxy}

\begin{proof}
\noindent\textbf{Distribution.}
The minimum of $D$ independent Geometric($p$) variables is Geometric($P_D$) where $P_D = 1-(1-p)^D$.
Thus $\E[T] = 1/P_D$ and $-\log P_D = \log \E[T]$.

\noindent\textbf{Jensen gap.}
By Jensen's inequality on the concave function $\log$:
$\E[\log T] \le \log \E[T] = -\log P_D$.
The gap $\Delta(p,D) = -\log P_D - \E[\log T] \approx (1-p)^D/2$ (second-order approximation using $\Var(T)/2\mu^2$ for geometric~$T$).

\noindent\textbf{Part (a): Fixed~$D$.}
$P_D = 1-(1-p)^D$ is strictly increasing in~$p$.
Thus $-\log P_D$ is strictly decreasing in~$p$.
By stochastic dominance, $\E[\log T]$ is also strictly decreasing in~$p$.
Rankings are identical.

\noindent\textbf{Part (b): Fixed model, varying~$D$.}
The true FOC is $\kappa = -\frac{d}{dD}(-\log P_D) + \frac{d}{dD}\Delta$.
Explicitly: $d\Delta/dD = -(1-p)^D \log(1-p)/2 < 0$ since $\log(1-p) < 0$ for $p \in (0,1)$.
Thus the true FOC's RHS is smaller than the proxy FOC's, satisfied at a smaller~$D$:
$\Dstar_{\mathrm{true}} \le \Dstar_{\mathrm{proxy}}$.
The difference is bounded by $|\Delta'|/|J''| = O((1-p)^{\Dstar})$, exponentially small.
\end{proof}

\subsection{Routing Triviality (Proposition~\ref{prop:routing-triviality})}
\label{app:routing}

\begin{proof}
Define cost-adjusted performance $\rho(s,M) = -\log p(s,M) + \log\kappa(M)$.

\noindent\textbf{Part (a): Sufficient condition.}
Let $M_c$ be the cheapest model.
For any $M'$ and any mode~$s$:
$\rho(s,M') - \rho(s,M_c) = (\log p(s,M_c) - \log p(s,M')) + (\log\kappa(M') - \log\kappa(M_c))$.
By the condition, $|\log p(s,M_c) - \log p(s,M')| < |\log\kappa(M') - \log\kappa(M_c)| = \log\kappa(M') - \log\kappa(M_c)$.
Thus $\rho(s,M') - \rho(s,M_c) > 0$ for all~$s$: $M_c$ wins everywhere.

\noindent\textbf{Part (c): Necessary condition.}
If routing is non-trivial, there exist $s_1, s_2$ and $M_1, M_2$ with $\rho(s_1,M_1) < \rho(s_1,M_2)$ and $\rho(s_2,M_1) > \rho(s_2,M_2)$.
This is a crossing in cost-adjusted log-competence---Ricardian comparative advantage.
Conversely, if no crossing exists, one model has the lowest $\rho$ on every mode, and routing is trivial.
\end{proof}
